\section{Các thách thức của bài toán}

Quá trình xây dựng một hệ thống RAG chuyên biệt và tin cậy cho tập dữ liệu về các bệnh lý về xương phải đối mặt với sáu nhóm thách thức cốt lõi, bao phủ toàn bộ luồng vận hành của hệ thống:

\subsection{Tính không đồng nhất, tính thời điểm và chất lượng của dữ liệu}
Kho tri thức y khoa được tổng hợp từ giáo trình, hướng dẫn lâm sàng, phác đồ và các tài liệu ở nhiều định dạng nên thường không đồng nhất về cấu trúc. Bảng, sơ đồ hoặc văn bản trích xuất từ PDF có thể bị mất quan hệ ngữ nghĩa; cùng một bệnh hoặc thuốc có thể mang nhiều tên gọi; các hướng dẫn giữa tổ chức, quốc gia và thời điểm ban hành có thể khác nhau hoặc mâu thuẫn. Nếu hệ thống không quản lý nguồn, phiên bản, phạm vi áp dụng và trạng thái hiệu lực, một đoạn văn bản lỗi thời vẫn có thể được truy xuất với điểm liên quan cao. Vì vậy, chất lượng phân đoạn và siêu dữ liệu của từng đoạn là điều kiện tiên quyết, không chỉ là bước tiền xử lý kỹ thuật.

\subsection{Giới hạn của bộ truy xuất trước rào cản thuật ngữ và nhiễu ngữ nghĩa}
Dữ liệu y khoa về các bệnh lý về xương sở hữu đặc thù với mật độ thuật ngữ chuyên ngành cao, bao gồm nhiều từ đồng nghĩa, từ mượn và danh từ viết tắt lâm sàng. Các thuật toán tìm kiếm theo từ khóa truyền thống như BM25 thường gặp hiện tượng bỏ sót tài liệu liên quan do sự bất đồng giữa từ ngữ người dùng truy vấn và thuật ngữ chính thống trong văn bản y học \cite{medrag2024}. Ngược lại, các mô hình nhúng biểu diễn không gian đa chiều tuy cải thiện khả năng tra cứu ngữ nghĩa nhưng lại dễ vấp phải rào cản nhiễu sau truy xuất \cite{xiong2024mirage}. Các mô hình này có xu hướng trích xuất các đoạn văn bản có độ tương đồng biểu diễn cao (như cùng đề cập đến một nhóm thuốc) nhưng lại thiếu hụt những thông tin lâm sàng mang tính quyết định (như chống chỉ định hoặc liều lượng áp dụng cho bệnh nhân có bệnh nền), dẫn đến sự sai lệch ngữ cảnh đầu vào cho mô hình ngôn ngữ lớn \cite{medtrustrag}.

\subsection{Đứt gãy ngữ cảnh và năng lực suy luận y khoa đa bước}
Một truy vấn y khoa trong thực tế lâm sàng thường đòi hỏi khả năng xử lý thông tin phức tạp và liên kết đa chiều. Để giải đáp một câu hỏi liên quan đến phác đồ điều trị kết hợp với tiền sử bệnh nền, hệ thống không thể chỉ dựa vào một đoạn văn bản đơn lẻ mà phải thực hiện tiến trình suy luận đa bước: xác định phác đồ điều trị chuẩn, sàng lọc danh sách chỉ định thuốc, và đối chiếu các tác dụng phụ đối với bệnh lý đi kèm \cite{imedrag}. Thách thức lớn nằm ở chỗ các kiến trúc RAG cơ bản thực hiện truy xuất đơn luồng thường làm đứt gãy mạch logic này. Việc chia nhỏ tài liệu thành các đoạn độc lập làm thất lạc thông tin liên kết, khiến mô hình ngôn ngữ lớn tiếp nhận ngữ cảnh thiếu hụt hoặc bị nhiễu loạn giữa tập tài liệu truy xuất có độ dài lớn \cite{xiong2024mirage}.

Trong hội thoại nhiều lượt, lịch sử trao đổi và bối cảnh bệnh nhân còn tích lũy thêm thông tin không liên quan, lỗi thời hoặc mâu thuẫn với câu hỏi hiện tại. Đưa nguyên văn toàn bộ lịch sử vào câu lệnh vừa chiếm cửa sổ ngữ cảnh vừa có thể làm mô hình ưu tiên sai bằng chứng. Hệ thống vì thế phải cô đọng lịch sử, giữ lại các dữ kiện lâm sàng cần thiết và viết lại câu hỏi nối tiếp thành truy vấn độc lập mà không làm mất ý định ban đầu.

\subsection{Xung đột tri thức và các dạng ảo giác y khoa}
Khi ngữ cảnh được trích xuất mang về các phác đồ y học mới nhưng lại mâu thuẫn với tri thức được huấn luyện trước trong tham số của mô hình ngôn ngữ lớn, hiện tượng xung đột tri thức sẽ xuất hiện. Sự mâu thuẫn này là nguyên nhân cốt lõi phát sinh bốn dạng ảo giác nguy hiểm trong tư vấn y tế \cite{medtrustrag, xiong2024mirage}:
\begin{enumerate}[label=(\arabic*)]
    \item \textbf{Suy luận lỗi \cite{medtrustrag}:} Mô hình kết nối sai logic giữa triệu chứng lâm sàng và phương pháp điều trị, dẫn đến việc đưa ra kết luận thiếu căn cứ xác thực từ tài liệu gốc.
    \item \textbf{Bỏ sót thông tin \cite{medtrustrag}:} Mô hình tóm tắt quá mức, làm bỏ sót các cảnh báo lâm sàng quan trọng hoặc các chống chỉ định y khoa nghiêm trọng.
    \item \textbf{Từ chối quá mức \cite{medtrustrag}:} Mô hình từ chối trả lời mặc dù dữ liệu ngữ cảnh đã cung cấp đầy đủ thông tin bằng chứng.
    \item \textbf{Gán sai nguồn trích dẫn \cite{medtrustrag}:} Mô hình đưa ra câu trả lời nhưng lại đính kèm trích dẫn tới đoạn văn bản không chứa thông tin chứng minh.
\end{enumerate}

\subsection{Sự khắt khe trong tiêu chuẩn an toàn và đánh giá lâm sàng}
Biên độ chấp nhận sai sót trong các ứng dụng y tế gần như bằng không. Câu trả lời của hệ thống không chỉ cần tuân thủ độ trung thực so với văn bản ngữ cảnh, mà còn phải đáp ứng các tiêu chuẩn lâm sàng nghiêm ngặt bao gồm độ chính xác y khoa và độ an toàn lâm sàng nhằm đảm bảo không gây hại cho người bệnh \cite{who2022musculoskeletal, medtrustrag}. Bên cạnh đó, do người sử dụng hệ thống bao gồm cả bệnh nhân và cộng đồng, câu trả lời còn phải đáp ứng yêu cầu về khả năng đọc hiểu (tránh sử dụng thuật ngữ quá hàn lâm gây khó hiểu) và sự đồng cảm trong giao tiếp y khoa \cite{ortho_rag_13}. Đây là những tiêu chuẩn đa chiều mà các bộ chỉ số đo lường tự động thông thường rất khó đánh giá một cách toàn diện.

Các độ đo tương đồng bề mặt như BLEU hoặc ROUGE không đủ để phản ánh rủi ro y khoa: một câu trả lời diễn đạt khác đáp án chuẩn vẫn có thể đúng, trong khi câu trả lời gần giống về từ ngữ vẫn có thể chứa một chi tiết nguy hiểm. Do đó, cần đánh giá riêng chất lượng truy xuất, độ chính xác y khoa, độ bám sát bằng chứng, tính đầy đủ, độ chính xác trích dẫn, khả năng từ chối hoặc chuyển chuyên gia và độ ổn định qua nhiều lượt hội thoại.

\subsection{Bảo vệ dữ liệu cá nhân và duy trì giám sát của con người}
Bối cảnh bệnh nhân và lịch sử hội thoại có thể chứa thông tin nhận dạng, tiền sử bệnh, thuốc đang dùng hoặc kết quả xét nghiệm nhạy cảm. Việc thu thập quá mức, lưu trữ không cần thiết hoặc đưa nguyên trạng các dữ liệu này vào dịch vụ mô hình có thể làm tăng nguy cơ lộ lọt thông tin. Hệ thống cần thực hiện giảm thiểu dữ liệu, ẩn danh và kiểm soát quyền sử dụng ngay từ khâu xây dựng kho tri thức đến suy diễn. Đồng thời, nguy cơ người dùng quá tin vào câu trả lời tự động đòi hỏi hệ thống phải biểu đạt độ không chắc chắn, nêu nguồn và duy trì cơ chế con người tham gia đối với các tình huống nguy cơ cao.
